\section{Summary and Future Directions}
\label{sec:rl_summary}

This chapter presented two contributions for enabling autonomous agents to adapt to novelty in open-world environments. The RAPid-Learn framework formalizes the integration of symbolic planning and reinforcement learning through the concept of integrated planning tasks, defines the executor discovery problem as the core challenge of novelty recovery, and introduces knowledge-guided exploration strategies that exploit domain knowledge to achieve sample-efficient adaptation. The neurosymbolic cognitive architecture embeds RAPid-Learn within a broader system that handles the full lifecycle of novelty---from detection through characterization to accommodation---by combining symbolic inference, neural inference, and reinforcement learning in a cascading recovery pipeline.

The experimental evaluations demonstrate three key findings. First, RAPid-Learn achieves one to two orders of magnitude faster adaptation than transfer learning baselines by decomposing the global recovery problem into focused subproblems defined by failed operators. Second, knowledge-guided exploration provides measurable benefits for complex novelties that require the agent to interact with novel entities or discover spatial relations not present in its symbolic model. Third, the integrated architecture achieves robust performance across 288 novelties spanning 10 categories, including concealed novelties not anticipated during system development, confirming that the combination of symbolic reasoning and learning-based recovery generalizes to a broad range of open-world changes.

Several limitations of the current approach point to important directions for future work.

\paragraph{Failure-driven adaptation.} As observed in \cref{ch:benchmarks}, the failure-driven nature of RAPid-Learn means that beneficial novelties that do not cause execution failures go undetected and unexploited. The chest novelty in the NovelGym evaluation (\cref{sec:ng_results}) illustrates this clearly: hybrid agents never discovered the shortcut because their plans continued to succeed. Extending the framework to proactively detect and exploit beneficial opportunities---for example, through periodic re-planning with updated models---would address this blind spot.

\paragraph{Fixed goal assumption.} The current framework assumes a fixed goal state throughout the adaptation process. In realistic open-world settings, goals themselves may need to be revised when novelty renders them unachievable or when new opportunities arise. Integrating goal reasoning~\citep{aha2018goalreasoning} with the executor discovery mechanism would enable more flexible responses to environmental changes.

\paragraph{Exploration cost in physical environments.} The reinforcement learning component of RAPid-Learn requires trial-and-error interaction with the environment, which is acceptable in simulation but potentially costly or dangerous in physical robotic settings. Reducing the number of exploratory interactions---through model-based reinforcement learning, sim-to-real transfer, or more informed exploration strategies---is essential for deploying these methods on real robots. This challenge motivates the work presented in \cref{ch:robot_interaction}, which develops reinforcement learning approaches for robotic manipulation that emphasize generalization across objects and embodiments, reducing the need for environment-specific exploration.

\paragraph{Symbolic-to-subsymbolic grounding.} The current architecture relies on a domain interface that provides pre-processed symbolic state descriptions. In real-world settings, the agent would need to ground these descriptions in raw sensory data---a challenging perception problem that interacts with novelty detection (e.g., visually recognizing a novel object and inferring its symbolic type). Tighter integration of the visual novelty detection capabilities with the symbolic reasoning components would enable end-to-end novelty handling from raw observations.

\paragraph{Coordination between adaptive components.} The systemic evaluation revealed that the interaction between the agent model and the executor learner can be counterproductive: the learner's exploration triggers false novelty detections from the agent model, leading to unnecessary exploration. Developing mechanisms for coordinating multiple adaptive components---for example, by suppressing novelty detection during active learning phases, or by sharing information about the source of environmental changes between components---is an important engineering and research challenge for complex novelty-handling architectures.